%! Author = matteomagnini
%! Date = 18/03/26

% Preamble
\documentclass[11pt]{article}

% Packages
\usepackage{amsmath}
\usepackage{amssymb}
\usepackage{my-acronyms}
\usepackage{hyperref}

\title{
    Intelligent System 2\\
    Tutorial week 5: Beliefs revision and update
}

\author{
    Matteo Magnini\\
    \href{mailto:matteo.magnini@uni.lu}{matteo.magnini@uni.lu}
}

\date{
    19 March, 2026
}

% Document
\begin{document}

    \maketitle

    \section{Exercises}\label{sec:exercises}

    \subsection{Tweety}

    \begin{itemize}
        \item It is the case that birds fly.
        \item Penguins are birds.
        \item Tweety is a penguin.
    \end{itemize}

    How do we revise our knowledge if we learn that Tweety does \emph{not} fly?

    \subsection{The Ring}

    Consider the following scenario:

    \begin{itemize}
        \item Gandalf does not know if the Ring has been used.
        \item Gandalf knows that if the ring has been used, then Sauron knows the location of the Ring.
        \item Gandalf also knows that if a Nazgûl knows the location of the Ring, then Sauron also knows the location of the Ring.
    \end{itemize}

    Frodo uses the Ring to escape the Nazgûl.

    What is the new knowledge base after this event, according to the belief update semantics?

\end{document}